% Standalone contribution for the E agent; the parent assigns a statement number.
\paragraph{The literal Corson--Obraztsov J\'onsson construction has large $H_2$.}
Let $\kappa$ be an uncountable cardinal admitting a J\'onsson algebra,
and choose the group $J_\kappa$ by Corson's proof, using the presentation
in Obraztsov's Theorem~A with $H=1$, $m=\infty$, and infinite cyclic
factors indexed by $\kappa$. Then
\[
                 H_2(J_\kappa;\mathbb Z)\cong\mathbb Z^{(\kappa)}.
\]
In particular this specific J\'onsson construction does not give an
acyclic group or a group with countable second homology.

\emph{Proof.}
Write $F=*_{\alpha<\kappa}\langle x_\alpha\rangle$ and
$J_\kappa=F/N$, where $N$ is normally generated by the defining
relators in Obraztsov's construction. Let $\mathcal R$ contain one
representative from each class under cyclic conjugacy and inversion;
retain the displayed representatives used below. His notion of
diagrammatic asphericity, stated on p.~268 and established on p.~285,
has the following sufficient consequence:
\begin{equation}\label{CD:EJonsson:relationcoinvariants}
                  N/[F,N]\cong\mathbb Z^{(\mathcal R)}.
\end{equation}
The convention is the one made explicit by Muranov, Section~6,
Definition~60: a spherical diagram splits into elementary two-face
spheres by diamond moves; his Definitions~33 and 38--40 describe
these moves and cancellations. Section~6 expressly identifies this
with Ol'shanskii's diagrammatic asphericity. In an oriented sphere,
diamond moves retain the oriented face boundary labels, and the
two faces of an elementary sphere have inverse boundary labels.
To see the consequence directly,
the displayed map is onto by normal generation. A relation in its
kernel can, after writing the relevant elements of $[F,N]$ as
commutators and expanding their $N$-entries into conjugates of
defining relators, be written as an identity product of conjugates
of relators. Its signed multiplicities are the originally proposed
relation. Draw the corresponding spherical identity diagram.
Removal of a cancelling pair of defining cells changes none of
these signed multiplicities. Repeated diagrammatic reductions end
with factor cells only; those have no defining-relator multiplicities.
Thus all the originally proposed multiplicities are zero. The
same argument works after expanding each letter of a cyclic factor
as a power of $x_\alpha$; factor cells fill in its circle and contribute
zero multiplicity. Every diagram involved is finite, so no
cardinality restriction enters this argument.

The five-term sequence for $1\to N\to F\to J_\kappa\to1$,
or equivalently Hopf's formula, now identifies
\begin{equation}\label{CD:EJonsson:Hopfkernel}
 H_2(J_\kappa;\mathbb Z)
 =\ker\!\left(\mathbb Z^{(\mathcal R)}
                 \xrightarrow{\partial}\mathbb Z^{(\kappa)}\right),
\end{equation}
where $\partial$ records ordinary exponent sums.

For each pair of distinct cyclic factors with indices $\alpha<\beta$,
the special periods listed on p.~277 include
$A=a[a,c]^t$, with nonidentity $a\in\langle x_\alpha\rangle$ and
$c\in\langle x_\beta\rangle$ and a permitted positive integer $t$.
Choose $a,c$ in the oriented halves used there; if the auxiliary
choice $a<b$ is needed, choose $a$ not maximal in its infinite
oriented half. Relation~(2.2) on that page includes, in particular,
the three relators with $k=6,7,8$ in the family
\[
 R_k=c^{-1}A^n
       \prod_{j=0}^{h-2}\bigl(c A^{n+30j+k}\bigr).
\]
Their exponent sums are
\[
 \partial R_k
 =\bigl(hn+15(h-2)(h-1)+(h-1)k\bigr)[a]+(h-2)[c].
\]
Consequently
\[
                    z_{\alpha\beta}=R_6-2R_7+R_8
\]
is a nonzero element of the kernel in
\eqref{CD:EJonsson:Hopfkernel}. The three relators are distinct:
their exponent sums in $x_\alpha$ are distinct, because $a$ is a
nonzero power of that generator. Their positive $x_\alpha$-coefficient
relative to $[a]$, together with the fixed nonzero $[c]$ coefficient,
also distinguishes their chosen orientations from inverse relators.

Partition the factor index set into $\kappa$ disjoint pairs.
The corresponding $z_{\alpha\beta}$ have disjoint supports in the
free abelian group on the relators, and are therefore independent.
Indeed every defining word in such a support uses exactly the
two factors of that pair. This proves that the rank of $H_2$ is
at least $\kappa$. There are at most $\kappa$ finite words over
the alphabet of the factors, hence at most $\kappa$ defining
relators. Formula~\eqref{CD:EJonsson:Hopfkernel} makes $H_2$ a
free abelian group of rank at most $\kappa$. Equality follows.
\hfill$\square$

\emph{Scope.}
This calculation uses the explicit presentation selected in
Corson's proof; it does not assert that every J\'onsson group has
large second homology. Nor does it identify Obraztsov's terminology
``diagrammatically aspherical'' with the existence of an actual
two-dimensional classifying space: his terminology also applies
to finite-exponent variants. The signed-relator argument and
Hopf's formula give the claimed obstruction without requiring
that additional geometric assertion.
